\section{Introduction}

Large language models (LLMs) now sit inside search engines, editors and
developer tools. They are reliable for knowledge that is common on the web,
but two kinds of information defeat them: information too large to read,
where accuracy drops even before the context limit is
reached~\cite{liu2024lost}, and information that is rare or private, such
as a company's internal systems, which the model saw seldom or never in
training~\cite{kandpal2023longtail}. In both cases models tend to fill the
gaps with plausible but unsupported claims~\cite{ji2023hallucination}.

Execution traces combine both problems. A kernel tracer such as
LTTng~\cite{desnoyers2006lttng} records every system call, scheduling
decision and interrupt; a 30-minute trace of the program we study holds
87 million events, far beyond any context window. And a developer's
question, such as why one thread is slow, concerns one specific program
that appears in no training corpus. Tools such as Trace
Compass~\cite{tracecompass} can answer it, but only for users who know
which analysis to open and how to read it.

Retrieval-augmented generation (RAG)~\cite{lewis2020rag} retrieves
relevant documents and adds them to the prompt. But documents explain what
a concept means or how a program was designed, not what happened in a
particular run: RAG can explain what a futex wait is, but not how long a
thread waited in one, or which thread woke it.

We present \sysname, a framework that answers plain-language questions
about kernel traces by grounding the LLM in computed evidence. Java
analyzers built on Trace Compass read the whole trace and produce measured
findings, such as how each thread's time divides between running, waiting
for CPU and being blocked, which thread woke which, and which thread
finishes each round last. RAG adds documentation, and the LLM explains the
findings rather than discovering them. Our contributions are:
\begin{itemize}
  \item \sysname, which grounds an LLM's answers in analyzer measurements
        of the whole trace, with RAG for documentation.
  \item A benchmark: LTTng traces of a multithreaded program with one slow
        thread at three lengths, and 13 questions with measured ground
        truth across five categories.
  \item A comparison using the same model: an LLM given raw trace excerpts
        answers 35\% correctly and an LLM with RAG 68\%, while \sysname
        answers 13 of 13 on every trace.
\end{itemize}
